\section*{Part III: Unit Economics: Accounting for 1 Repeatable Unit}

\subsection*{18. Contribution Margin Is the 1st Economic Truth}

For a unit with net realized revenue $P$ and incremental serving cost $c$, unit contribution is $P - c$. The contribution margin ratio is $(P - c)/P$ when $P > 0$. The contribution helps cover fixed costs, reinvestment, financing, and eventual profit. A product with positive contribution can still lose money if volume is too low; a negative-contribution product usually becomes more expensive to operate as volume increases.

But a founder must state which costs are included. Is customer acquisition cost (CAC) counted in the unit? Are refunds and payment fees deducted from revenue or treated as expenses? Is skilled teaching labor fixed per cohort or variable per learner? Different definitions can be valid for different decisions, but the figures cannot be compared without reconciliation.

Use at least 2 layers:
\begin{enumerate}
  \item \textbf{Delivery contribution:} revenue minus directly variable serving costs.
  \item \textbf{Customer-level contribution:} delivery contribution minus acquisition and directly variable sales costs.
\end{enumerate}
Then deduct fixed cohort costs and shared overhead. This makes the economics readable without falsely classifying every expense as a marginal cost.

\subsection*{19. Gross Margin, Contribution, and Operating Profit Are Not Synonyms}

Gross margin generally relates recognized revenue to cost of goods or services sold under the firm's accounting policy. Contribution margin is a managerial measure that groups costs by behavior relevant to an operating decision. Operating profit deducts operating expenses, including relevant fixed overhead and depreciation, according to the financial reporting framework. 1 firm's ``gross margin'' may not be comparable with another's if labor and support are classified differently.

For a software business, infrastructure is often included in cost of revenue, while many sales and product-development costs appear below gross profit. A consulting business's professional labor is a major cost of service. An educational provider should not exclude every teacher from delivery cost just because the teacher is also the founder. \textbf{A high gross margin does not imply healthy net income, and a high contribution percentage can coexist with a dangerous cash cycle.}

\subsection*{20. A Complete 1-Cohort Contribution Statement}

Consider an illustrative cohort of 18 learners paying ₹35,000 each. Assume no discounts or refunds for this simplified model. Revenue is $18 \times ₹35,000 = ₹6,30,000$. Per learner, directly variable processing, materials, and feedback cost ₹4,000, totaling ₹72,000. Acquisition and sales expenditure attributable to the 18 paid enrollments is ₹90,000. Fixed cohort teaching and preparation costs ₹2,10,000. The resulting contribution after cohort-specific fixed costs is \textbf{₹2,58,000} before shared institutional overhead, tax, and financing (Table~\ref{tab:cohort-contribution-statement}).

\begin{table}[htbp]
\centering
\small
\begin{tabularx}{\columnwidth}{Xr}
\toprule
\textbf{Item} & \textbf{Cohort Amount} \\
\midrule
Recognized cohort revenue ($18 \times ₹35,000$) & ₹6,30,000 \\
Variable learner delivery ($18 \times ₹4,000$) & $-$₹72,000 \\
\midrule
\textbf{Delivery contribution before cohort-fixed costs} & \textbf{₹5,58,000} \\
Acquisition and directly attributable sales & $-$₹90,000 \\
Fixed teaching, design, and cohort operations & $-$₹2,10,000 \\
\midrule
\textbf{Contribution after cohort-specific costs} & \textbf{₹2,58,000} \\
\bottomrule
\end{tabularx}
\caption{Illustrative 18-learner cohort economic contribution statement.}
\label{tab:cohort-contribution-statement}
\end{table}

This is deliberately not called net profit. Nor does the existence of ₹2,58,000 on paper establish that the founder has the capacity to teach the program twice. Examine actual contact hours, preparation, support, administrative time, and whether the same price and acquisition cost recur with the next cohort.

\subsection*{21. Cash Contribution Versus Normalized Economic Contribution}

In the earliest stage a founder may deliver teaching without a salary or hire. The cash ledger can therefore show attractive surplus. But a model that collapses after replacing the founder with a competent instructor has not demonstrated organizational independence.

Maintain 2 views:
\begin{itemize}
  \item \textbf{Cash contribution:} includes actual cash payments and receipts needed to run the pilot.
  \item \textbf{Normalized contribution:} includes a defensible economic cost of the founder's delivery and management labor, even if unpaid.
\end{itemize}
Do not automatically book an invented salary as a statutory expense: this is a managerial adjustment, not an instruction to falsify accounts. The difference is strategically valuable because it shows how much apparent profit is actually uncompensated founder labor.
